%======================================================================
\section{Further derivations}\label{app:derivations}
%======================================================================

\subsection{Isometry for the compound Poisson integral}

Let $\tilde L_t=L_t-\lambda_Jm_Jt$ be the compensated compound Poisson process. For a
deterministic square-integrable $f$, $\E\int_0^tf(u)\dd\tilde L_u=0$ and
\[
  \E\Bigl[\Bigl(\int_0^tf(u)\dd\tilde L_u\Bigr)^2\Bigr]=\lambda_J\E[J^2]\int_0^tf(u)^2\dd u .
\]
\emph{Proof.} For a step function $f=\sum_kf_k\1_{(u_k,u_{k+1}]}$ the increments of $\tilde L$
over disjoint intervals are independent with mean zero and variance
$\lambda_J\E[J^2](u_{k+1}-u_k)$, so the second moment is
$\lambda_J\E[J^2]\sum_kf_k^2(u_{k+1}-u_k)$. The general case follows by density of step
functions in $L^2$. Writing $\int f\dd L=\int f\dd\tilde L+\lambda_Jm_J\int f\dd u$ gives the mean
and variance used in Proposition~\ref{prop:cet1}, with $f(u)=e^{-\kappa(t-u)}$.

\subsection{Transition of the CET1 process between observation dates}

Conditionally on $C_s$, \eqref{eq:cet1sol} applied on $[s,t]$ gives
\[
  C_t=\theta+(C_s-\theta)e^{-\kappa(t-s)}+\sigma\int_s^te^{-\kappa(t-u)}\dd W_u-\int_s^te^{-\kappa(t-u)}\dd L_u .
\]
Without jumps the transition is Gaussian with mean $m=\theta+(C_s-\theta)e^{-\kappa\Delta}$ and
variance $\sigma^2(1-e^{-2\kappa\Delta})/(2\kappa)$, $\Delta=t-s$, which is \eqref{eq:outrans}.
With jumps it is a Gaussian mixture: conditionally on $k$ jumps at times
$u_1,\dots,u_k$, the mean is shifted by $-\sum_iJ_ie^{-\kappa(t-u_i)}$. The likelihood used in
Section~\ref{sec:calibration} ignores this term. Over a quarter the probability of a jump is
$1-e^{-0.12/4}\approx3\%$, so the Gaussian likelihood treats rare large moves as diffusion,
which biases $\hat\sigma$ upwards slightly and, through the profile likelihood, $\hat\kappa$ in
an ambiguous direction; the 2023Q1 exclusion removes the one large positive move.

\subsection{Profile likelihood}

Differentiating \eqref{eq:loglik} with respect to $\sigma^2$ gives
$\partial\ell/\partial\sigma^2=-\frac{n}{2\sigma^2}+\frac1{2\sigma^4}\sum_k(C_{k+1}-m_k)^2/w_k$,
which vanishes at \eqref{eq:sigmahat}. Substituting back,
\[
  \ell_p(\kappa)=\ell\bigl(\kappa,\hat\sigma(\kappa)\bigr)
  =-\frac n2\Bigl[1+\log\bigl(2\pi\hat\sigma^2(\kappa)\bigr)\Bigr]-\frac12\sum_k\log w_k(\kappa),
\]
a one-dimensional function of $\kappa$. The likelihood ratio statistic for a fixed $\kappa_0$
is $2[\ell_p(\hat\kappa)-\ell_p(\kappa_0)]$, asymptotically $\chi^2_1$.

\subsection{Yield to call and the delta method}

Let $y(V)$ solve $V=\sum_jc_j(1+y/2)^{-2t_j}$ for the bullet cash flows $c_j$ to $T_1$. By the
implicit function theorem
\[
  \frac{\dd y}{\dd V}=-\Bigl[\sum_jc_jt_j(1+y/2)^{-2t_j-1}\Bigr]^{-1}=-\frac{1}{D_{\mathrm{mod}}V},
\]
where $D_{\mathrm{mod}}$ is the modified duration to call. For a zero-recovery bullet to
$T_1$ with a constant hazard $h$, the price is the risk-free price discounted at $r+h$, so the
continuously compounded yield rises one-for-one with the hazard; in the base attribution the
hazard block of $219$bp of hazard produces $231$bp of semi-annual yield-to-call spread, a ratio
of $1.05$ that reflects compounding and the state dependence of the driver. The price
sensitivity $|V_0'|\approx0.08$ points per basis point of hazard is larger than a bullet's,
because it includes the exposure of extended cash flows beyond $T_1$. The delta method for the
implied hazard, \eqref{eq:sehazard}, follows from the same implicit function argument applied
to $V_0(\lambda)=P$.

\subsection{The credit triangle for the CDS reference points}

For a CDS with constant hazard $h$, recovery $R$ and continuous premium $s$, the protection
leg is $(1-R)\int_0^Th e^{-(r+h)u}\dd u$ and the premium leg $s\int_0^Te^{-(r+h)u}\dd u$, so the
par spread is $s=h(1-R)$ and $h=s/(1-R)$. The CDS rows of Table~\ref{tab:ladder} use this
approximation; it ignores the term structure of the hazard, accrued premium on default and
the difference between the standard coupon and the par spread, all of which are small at
these spread levels.

\subsection{Relation of the nested value to the Owen value}

For a game $v$ on $\mathcal N$ with a partition into unions $\{\mathcal H,\mathcal K\}$, the
Owen value of player $j\in\mathcal H$ averages its marginal contribution over all orders in
which the members of each union are contiguous, with both orders of the unions equally
likely \citep{Owen1977}. The nested value of Section~\ref{sec:attribution} restricts these
orders to those in which $\mathcal H$ comes first. It is therefore the Owen value of the game
in which the union order is fixed, and it inherits efficiency, symmetry within unions and
the null-player property. The full Owen value would average the nested value with the
attribution in which the contract block comes first, which values extension on a riskless
bond and is excluded for the reason given in Section~\ref{sec:robust}.
